%! Author = denis
%! Date = 11.06.2026

% Preamble
\documentclass[11pt]{article}

% Packages
\usepackage{amsmath}
\usepackage[margin=2.5cm]{geometry}

\title{Evaluation of the TMS Software Requirements Specification}
\author{Denis Schatzmann}
\date{\today}

% Document
\begin{document}

\maketitle

All the answers are about the final version 1.0 of the document, from January 2003.

\section{Who created the final version?}
Version 1.0 (31 December 2002) was done by Osamah Yacoub and Ahmad Arrabi. The
company that wrote the document is AlliedSoft, for a USAID program in Jordan.

\section{What is the intended use?}
The document writes down everything the system has to do, so the client can
agree to it and the team can build and test it. The system itself is a website
that is used to manage the training and testing of government employees in
Jordan.

\section{Who are the intended users?}
The users are listed in section 4.2. There is a system administrator, two kinds
of other administrators, a training coordinator, and a normal employee. Each of
them can do different things in the system.

\section{What is the standard environment?}
It is a website that runs over the internet or a company network. You open it in
Internet Explorer (version 5.0 or newer) at a screen size of 800 by 600, and it
works in English or Arabic. The client provides the computers and software it
runs on.

\section{Evaluation of a functional requirement}
I picked requirement FR19.4. It says that a response can have a file attached to
it, and the file can be at most 2 MB.
\begin{itemize}
  \item Complete: only partly. It does not say which file types are allowed.
  \item Accurate: only partly. It is not clear if you can attach one file or
        several, because other parts of the document say different things.
  \item Unambiguous: no. It is not clear what kind of file is meant or how many
        files you can add.
  \item Traceable: yes. The requirement has an ID and points to the matching use
        cases.
  \item Testable: yes. You can just try a file under 2 MB (should work) and one
        over 2 MB (should be blocked).
\end{itemize}

\section{Is there a requirement regarding data security?}
Yes. Section 5.2 says that different users get different access rights. There are
four levels, and you have to log in to get your level. It only talks about
logging in and what each user is allowed to do. It does not say anything about
things like encryption or backups.

\section{Is there a requirement how the software should be designed?}
Yes. One requirement says the system should be built in three layers so it can
work together with other government applications later. Another one says the team
has to write a separate document just for the software design.

\end{document}
